% ============================================================
% DOCUMENTCLASS
% ============================================================
\documentclass[12pt,oneside]{book}

% ============================================================
% METADATOS
% ============================================================
\newcommand{\universidad}{Universidad de Buenos Aires}
\newcommand{\facultad}{Facultad de Derecho}
\newcommand{\carrera}{Clínica Jurídica}
\newcommand{\materia}{Clínica Jurídica · Práctica en Derecho de Familia}
\newcommand{\materiacorta}{Clínica Jurídica}
\newcommand{\titulo}{Apuntes de Clínica Jurídica: Práctica en Derecho de Familia}
\newcommand{\autor}{Isaac Edgar Camacho Ocampo}
\newcommand{\anio}{2026}

% ============================================================
% PAQUETES
% ============================================================
\usepackage[utf8]{inputenc}
\usepackage[T1]{fontenc}
\usepackage[spanish,es-nodecimaldot]{babel}

\usepackage[
    top=2.5cm,
    bottom=2.5cm,
    left=3cm,
    right=2.5cm,
    headheight=15pt
]{geometry}

\usepackage{amsmath}
\usepackage{amssymb}
\usepackage{mathtools}

\usepackage{graphicx}
\usepackage{float}
\usepackage{caption}
\usepackage{subcaption}

\usepackage{booktabs}
\usepackage{array}
\usepackage{longtable}
\usepackage{tabularx}

\usepackage{enumitem}

\usepackage{xcolor}
\usepackage[most]{tcolorbox}

\usepackage{fancyhdr}
\usepackage{ragged2e}

\usepackage[
    colorlinks=true,
    linkcolor=blue!50!black,
    citecolor=green!40!black,
    urlcolor=blue!60!black,
    pdftitle={\titulo},
    pdfauthor={\autor},
    pdfsubject={Apuntes universitarios},
    bookmarks=true
]{hyperref}

% ============================================================
% CONFIGURACIÓN
% ============================================================
\graphicspath{
    {./img/}
    {./graficos/}
    {./figuras/}
}

\setcounter{tocdepth}{2}

\setlength{\parindent}{0pt}
\setlength{\parskip}{0.5em}

\setlist[itemize]{
    topsep=0.3em,
    itemsep=0.2em
}

\setlist[enumerate]{
    topsep=0.3em,
    itemsep=0.2em
}

\clubpenalty=10000
\widowpenalty=10000

\pagestyle{fancy}
\fancyhf{}
\fancyhead[L]{\nouppercase{\leftmark}}
\fancyhead[R]{\materiacorta}
\fancyfoot[C]{\thepage}
\renewcommand{\headrulewidth}{0.4pt}
\renewcommand{\footrulewidth}{0pt}

\fancypagestyle{plain}{
    \fancyhf{}
    \fancyfoot[C]{\thepage}
    \renewcommand{\headrulewidth}{0pt}
}

% ============================================================
% COMANDOS
% ============================================================
\newcommand{\abs}[1]{\left|#1\right|}
\newcommand{\norm}[1]{\left\lVert#1\right\rVert}
\newcommand{\vect}[1]{\mathbf{#1}}
\newcommand{\deriv}[2]{\frac{d#1}{d#2}}
\newcommand{\pderiv}[2]{\frac{\partial#1}{\partial#2}}

% ============================================================
% ENTORNOS / CAJAS
% ============================================================
\newtcolorbox{definition}[1]{
    enhanced,
    breakable,
    title={Definición: #1},
    fonttitle=\bfseries,
    colback=blue!3,
    colframe=blue!50!black,
    boxrule=0.8pt,
    arc=2mm
}

\newtcolorbox{important}{
    enhanced,
    breakable,
    title={Importante},
    fonttitle=\bfseries,
    colback=yellow!8,
    colframe=orange!70!black,
    boxrule=0.8pt,
    arc=2mm
}

\newtcolorbox{example}[1]{
    enhanced,
    breakable,
    title={Ejemplo: #1},
    fonttitle=\bfseries,
    colback=green!4,
    colframe=green!40!black,
    boxrule=0.8pt,
    arc=2mm
}

\newtcolorbox{formula}[1]{
    enhanced,
    breakable,
    title={Fórmula / Regla: #1},
    fonttitle=\bfseries,
    colback=purple!4,
    colframe=purple!50!black,
    boxrule=0.8pt,
    arc=2mm
}

\newtcolorbox{exercise}[1]{
    enhanced,
    breakable,
    title={Ejercicio: #1},
    fonttitle=\bfseries,
    colback=gray!5,
    colframe=gray!60!black,
    boxrule=0.8pt,
    arc=2mm
}

\newtcolorbox{note}{
    enhanced,
    breakable,
    title={Nota},
    fonttitle=\bfseries,
    colback=white,
    colframe=gray!50,
    boxrule=0.6pt,
    arc=2mm
}

% ============================================================
% CONFIGURACIÓN FINAL
% ============================================================
\renewcommand{\arraystretch}{1.25}

% ============================================================
% DOCUMENT
% ============================================================
\begin{document}

% ------------------------------------------------------------
% PORTADA
% ------------------------------------------------------------
\begin{titlepage}

\thispagestyle{empty}

\begin{center}

\vspace*{1cm}

{\large \textsc{\universidad}}

\vspace{0.2cm}

{\large \textsc{\facultad}}

\vspace{3cm}

{\Huge \textbf{\materiacorta}}

\vspace{1cm}

{\LARGE Práctica en Derecho de Familia}

\vspace{0.8cm}

\textit{Comprender los conceptos es el primer paso para convertir el estudio en conocimiento.}

\vspace{1.2cm}

\rule{0.70\textwidth}{0.5pt}

\vspace{0.5cm}

{\large Apuntes de estudio}

\vspace{0.5cm}

\rule{0.70\textwidth}{0.5pt}

\vfill

{\large \autor}

\vspace{0.3cm}

{\large \anio}

\end{center}

\vfill

\begin{flushleft}
\textit{Este es un modesto aporte a los estudiantes de ``\materia'' no pretende sustituir el material oficial de la materia.}
\end{flushleft}

\end{titlepage}

% ------------------------------------------------------------
% ÍNDICE
% ------------------------------------------------------------
\tableofcontents

% ------------------------------------------------------------
% CONTENIDO
% ------------------------------------------------------------

\chapter*{Presentación de la clase}
\addcontentsline{toc}{chapter}{Presentación de la clase}
\markboth{Presentación de la clase}{Presentación de la clase}

Los apuntes corresponden a una clase de práctica de la Clínica Jurídica en Derecho de Familia, de octubre de 2026 (según las referencias de la clase), desarrollada con casos reales de personas asistidas y registrada en dos grabaciones consecutivas. El tema central es \textbf{cómo se confecciona el informe de caso} y \textbf{cómo se conduce la relación con el asistido}.

La clase tuvo dos partes:

\begin{itemize}
    \item En la primera, el docente fijó el \textbf{estándar mínimo del informe de caso}: datos completos del juzgado, medidas provisorias con su vencimiento, deuda y pagos, composición del grupo familiar y respuesta de los oficios.
    \item En la segunda, el equipo repasó las novedades de varios casos (alimentos a un demandado sin recursos, denuncias paralizadas, una mediación cerrada por domicilio desconocido y la estrategia ante el incumplimiento) y atendió a una consultante. En esa reunión se le explicaron los \textbf{alimentos provisorios}, se le marcó un \textbf{límite ético} (no contactar a la abogada contraria) y se la preparó para la entrevista del \textbf{CIF}.
\end{itemize}

\begin{note}
Por reserva profesional, en el material de origen los nombres de las personas asistidas, de la contraparte, de sus abogadas y los domicilios fueron reemplazados por su rol. Lo que figura entre corchetes en el material original son agregados para dar sentido a frases incompletas; los fragmentos ininteligibles fueron omitidos.
\end{note}

\section*{Utilidad profesional de lo trabajado}
\addcontentsline{toc}{section}{Utilidad profesional de lo trabajado}

El informe de caso es el \textbf{núcleo del oficio}: un informe bien hecho permite que cualquier colega retome el caso sin ir al expediente, protege al estudio ante una ausencia o reemplazo y evita errores en lo que se reclama. Además, enseña a manejar al cliente (decirle lo que no quiere oír, ponerle límites, prepararlo para una pericia) y a cumplir las reglas de conducta profesional: comunicarse sólo entre abogados y preservar la distancia emocional.

Esta utilidad alcanza también a quien no ejerza: se aprende a \textbf{documentar con orden}, a \textbf{entender un proceso desde afuera sin involucrarse emocionalmente} y a \textbf{comunicar malas noticias con empatía y claridad}. En palabras del docente, aun quien no vaya a ejercer se beneficia de tener «la costumbre de saber qué es lo que están haciendo con su trabajo, conocer a su cliente, conocer la situación».

\begin{example}{los bloques de la clase como la preparación de una valija}
Los apuntes originales proponen una analogía cotidiana: armar la valija para el viaje de un familiar que viaja solo. El \textbf{informe de caso} (Bloque I) es la lista de lo que hay que llevar; las \textbf{novedades y estrategias} (Bloque II) son los imprevistos del viaje; la \textbf{organización del equipo} (Bloque III) es quién carga qué y a qué hora salen; la \textbf{reunión con la consultante} (Bloque IV) es hablar con el viajero para que no improvise; y la \textbf{familia ampliada} (Bloque V) son los parientes que ayudan pero no conducen. Si falta un solo eslabón, el caso se complica.
\end{example}

% ============================================================
\chapter[Informe de caso]{Bloque I · El estándar del informe de caso}
% ============================================================

\section{Datos del juzgado y del expediente}

El informe debe consignar los \textbf{datos completos del juzgado}: dirección, piso, secretaría, nombre del juez y del secretario. El docente señaló que estos datos deben saberse «de cabeza» y que casi ningún informe los trae.

% TODO: contenido incompleto o ambiguo en las notas originales (el inicio de la intervención del docente en la primera grabación es ininteligible).

\section{Medidas provisorias y sus vencimientos}

Cuando existan medidas provisorias, el informe debe indicar \textbf{cuántas} hay, \textbf{de qué tipo} son (alimentos, restricciones o cualquier otra) y la \textbf{fecha de vencimiento} de cada medida. El docente calificó este dato como «fundamental» y pidió que todos lo anotaran.

\begin{important}
En el informe de caso deben figurar: \textbf{cuántas} medidas provisorias hay, \textbf{de qué tipo} son (alimentos, restricciones, etc.) y \textbf{cuándo vencen}.
\end{important}

\section{Cuantificación del reclamo: deuda y pagos}

Si se va a reclamar dinero, el equipo necesita saber \textbf{cuánto se debe}. No se exige una liquidación prolija y presentable en el expediente, pero sí una \textbf{aproximación}: sin ella no se sabe qué monto corresponde pedir. También debe constar \textbf{si el obligado pagó algo}. Según el docente, estos datos tienen que estar en el informe «sí o sí»; un informe que no los contiene no se lee.

\section{Composición del grupo familiar y datos del entorno}

El informe debe contener, como mínimo:

\begin{itemize}
    \item cuántos hijos son y qué edad tienen;
    \item si los hijos estudian y dónde;
    \item si tienen obra social;
    \item si trabaja otra persona;
    \item si trabaja la otra parte.
\end{itemize}

Estos datos son especialmente relevantes en materia de alimentos. El docente manifestó su malestar ante informes que omiten estos datos, porque indican que \textbf{no se leyó la causa}. Un informe bien hecho permite al docente leerlo sin acudir a la causa: \textbf{si debe ir al expediente, es porque el informe está mal hecho}.

\begin{important}
El informe no consiste en cortar y pegar el contenido del drive, sino en \textbf{trabajar} y consignar lo que el equipo considere importante además de los datos básicos. Puede haber errores de apreciación (para eso está el docente), pero las cosas básicas no pueden faltar: los datos de cabeza del juzgado, la existencia de hijos menores y si trabaja la otra parte.
\end{important}

\begin{example}{oficios respondidos pero no registrados}
A veces las entidades respondieron los oficios y el equipo docente no lo sabe: no se informa lo que respondió Mercado Libre ni lo que respondió el lugar de trabajo. Esa información debe ser aportada por los estudiantes en el informe.
\end{example}

\section{El informe como trabajo en equipo y como práctica profesional}

El docente explicó que depende del equipo estudiantil: tiene sus propios trabajos y estudios y no puede ser tan eficiente como quisiera. Los estudiantes son quienes trabajan directamente con el consultante y saben qué decir y qué no; si saben que, por ejemplo, la hermana trabaja o que el demandado trabaja en una empresa con sueldo, ese dato debe ser aportado por ellos, porque el docente no lo averiguará si no lo escucha casualmente. En ese esquema, \textbf{«nosotros vamos y ponemos la cara, ustedes [ponen la información]»}.

Además, el informe constituye una \textbf{buena práctica profesional}:

\begin{itemize}
    \item Ante un cúmulo de trabajo, o si la vida personal no permite seguir con atención cada caso, el informe es un «machete» perfecto: el \textbf{mapa completo de la situación} de la persona que consultó.
    \item Si debe intervenir un socio, un reemplazante o un colega de emergencia, ya tiene una «foto» del caso para continuar el trabajo.
    \item Aunque no se ejerza, sirve para acostumbrarse a conocer el trabajo, al cliente y la situación. El docente agregó que además se les va a exigir.
\end{itemize}

\begin{important}
Un buen informe permite que cualquier colega \textbf{retome el caso sin ir al expediente}. Datos mínimos obligatorios: juzgado completo, medidas provisorias y vencimientos, deuda y pagos, hijos (cantidad, edades, estudios, obra social), situación laboral de las partes y respuestas de oficios.
\end{important}

\subsection*{Cuadro Cornell del Bloque I}

\begin{table}[H]
\centering
\begin{tabularx}{\textwidth}{
    >{\raggedright\arraybackslash}p{0.26\textwidth}
    >{\raggedright\arraybackslash}X
}
\toprule
\textbf{Preguntas / palabras clave} &
\textbf{Notas / respuestas} \\
\midrule

\textbf{¿Qué datos mínimos debe tener todo informe de caso?} &
Juzgado completo (dirección, piso, secretaría, juez y secretario), medidas provisorias con su vencimiento, deuda y pagos, composición del grupo familiar (hijos, edades, estudios, obra social), situación laboral de las partes y respuestas de los oficios. \\

\textbf{¿Cuál es la prueba de que un informe está mal hecho?} &
Que el docente tenga que ir al expediente para entender el caso. Un buen informe permite leerlo sin acudir a la causa. \\

\textbf{¿Por qué no basta con «cortar y pegar» del drive?} &
Porque el informe exige trabajo propio: seleccionar y destacar lo importante luego de los datos básicos. Pegar sin leer la causa produce un informe sin utilidad. \\

\textbf{¿Qué debe saber el equipo sobre la deuda alimentaria?} &
Cuánto se debe (liquidación aproximada) y si el obligado pagó algo; sin eso no se sabe qué reclamar. \\

\textbf{¿Qué ventaja profesional tiene un buen informe?} &
Funciona como «mapa» del caso: permite reemplazar o reforzar al abogado a cargo y, aunque no se ejerza, enseña a conocer el trabajo y al cliente. \\

\midrule

\multicolumn{2}{p{\dimexpr\textwidth-4\tabcolsep\relax}}{
\textbf{Resumen:}
El informe de caso debe ser completo y elaborado con trabajo propio, de modo que cualquier colega pueda retomar el caso sin acudir al expediente. Debe incluir los datos del juzgado, las medidas provisorias y sus vencimientos, la deuda y los pagos, el grupo familiar y las respuestas de los oficios.
} \\

\bottomrule
\end{tabularx}
\end{table}

% ============================================================
\chapter[Novedades y estrategia]{Bloque II · Novedades de los casos y estrategia}
% ============================================================

\section{Avisos de la clase y novedades con la contraparte}

\subsection{Avisos generales}

\begin{itemize}
    \item \textbf{Fecha de parcial:} el día 22. El parcial toma las causas propias de cada estudiante.
    \item \textbf{Mascota en una de las causas:} la entrega del gato continuaba pendiente. Se informó que se habían fijado alimentos y que la entrega todavía no se había realizado; la contraparte había quedado en comunicarse con la abogada para definir la forma de entrega. El docente consideró que lo actuado estuvo bien «porque no es ético», pero que había que entregar el gato.
\end{itemize}

% TODO: contenido incompleto o ambiguo en las notas originales (la frase del docente sobre la ética y la entrega del gato no permite precisar a qué conducta se refiere).

\subsection{Informe del cuerpo interdisciplinario}

Se recibieron informes del CIF: indicaron que la niña será evaluada nuevamente y que se convocó al equipo. En el informe del CIF, a favor de la posición del equipo, se señaló que la \textbf{multiplicidad de organismos que realizan el mismo informe es perjudicial} porque \textbf{revictimiza a la niña}; por lo tanto, se propone que lo realice el Ministerio Público.

\subsection{Plazos de contestación de demanda}

La demanda debía contestarse dentro del plazo, que vencía el jueves. Se agregaron los testigos disponibles y el escrito ya había sido visto, con correcciones menores. Resta realizar el petitorio. El escrito incluye también el tema del Instagram.

\section{Alimentos fijados a un demandado sin conocimiento ni recursos}

Un demandado, adulto mayor y veterano de guerra, manifestó estar muy mal y no entender qué es el «alimento»: cree que consiste en llevar comida a la casa. No maneja bancos ni sabe cómo cumplir. Quiere ir a la casa y entregar lo que corresponda. Su único ingreso, según lo informado, es el puesto que tenía en la casa y una pensión por ser veterano; se estima que sobrevive con dificultad. Se planteó el interrogante de cómo lograr que cumpla los alimentos y la necesidad de notificarlo.

\begin{example}{la cuota alimentaria que el obligado no comprende}
Un demandado adulto mayor, con ingresos mínimos (una pensión), recibe una cuota alimentaria fijada por el juzgado y \textbf{entiende «alimentos» como llevar comida a la casa}; no maneja bancos ni transferencias. La sentencia existe, pero el cumplimiento exige \textbf{notificarlo bien, explicarle la obligación y definir una forma de pago acreditable}.
\end{example}

\section{Denuncias de violencia paralizadas por falta de impulso}

Respecto de las causas por violencia, la consultante las presentó «muy por arriba» y quedaron paralizadas porque no se impulsaron. No se trata de denuncias de oficio, sino únicamente de denuncias de ella contra él. Se mencionaron las causas de 2023 y de 2020; esta última (la más reciente según lo dicho) refiere a maltrato a la hija: la madre denunció que el padre retiraba a la niña del colegio y que ella volvía con moretones. La denuncia se hizo en la OVD y hubo dos o tres actuaciones, pero «no dice nada» porque no se impulsó. Los estudiantes que expusieron admitieron no haber leído el detalle de esas causas.

% TODO: contenido incompleto o ambiguo en las notas originales (el orden entre las causas de 2020 y 2023 y su estado exacto no se precisa).

\begin{important}
Una denuncia sin impulso queda paralizada y \textbf{no produce prueba ni medidas}. Para la estrategia del caso importa lo que efectivamente se tramitó, no lo que se denunció. Por eso el informe debe indicar el \textbf{estado real} de cada causa vinculada, y la lectura de esas causas es un deber del equipo.
\end{important}

\section{Mediación fallida: domicilio desconocido}

El centro de mediación informó que la parte requerida no estaba notificada porque el correo respondió «domicilio desconocido». El docente advirtió que el domicilio consignado como real era un \textbf{domicilio laboral} (un comercio), sin piso ni departamento. La decisión fue \textbf{cerrar la mediación}: si más adelante se conoce otro domicilio, se reabre. El centro señaló que, si hubiera otro domicilio, la mediación no serviría, y remitió el acta de cierre para su firma.

\begin{example}{cierre de una mediación por falta de domicilio real}
La parte requerida no pudo ser notificada porque el único domicilio conocido era \textbf{laboral} y el correo informó «domicilio desconocido». Decisión del docente: \textbf{cerrar la mediación} y, si luego se conoce un domicilio real, \textbf{reabrirla}. Luego se firma el acta de cierre.
\end{example}

\section{Inicio de la demanda y estrategia frente al incumplimiento}

El docente anunció que se iniciaría la demanda; la contraparte cuenta también con una abogada. Se estimó un plazo de dos o tres semanas para armarla correctamente.

Sobre el incumplimiento:

\begin{itemize}
    \item Ya existen \textbf{alimentos provisorios} fijados y el obligado \textbf{no paga}. Se ejecutará; el obligado no tiene cuenta bancaria.
    \item La deuda existe, pero no puede asegurarse si se cobrará en un mes o en tres. Se trabaja para poder ejecutar aunque sea una suma como 200.000 pesos.
    \item El docente se refirió a una estrategia de la otra parte: empezar a pagar 450 [mil] y no pagar más, para decir «no tengo» y presionar a que el alimentado prefiera aceptar esa suma antes que seguir peleando por 600 o 700, o la cifra que fuere. Calificó esto como una \textbf{estrategia psicológica}; frente a ella, el equipo decidió \textbf{litigar hasta el final}.
\end{itemize}

% TODO: contenido incompleto o ambiguo en las notas originales (las cifras 450, 600 y 700 se mencionan sin unidad ni contexto preciso).

\subsection{Herramientas de ejecución mencionadas}

En los provisorios, \textbf{si hay tres cuotas impagas}, puede pedirse la \textbf{inscripción en el registro de deudores alimentarios}, con la consecuencia de que el deudor no podrá renovar el registro de conducir.

\subsection*{Cuadro Cornell del Bloque II}

\begin{table}[H]
\centering
\begin{tabularx}{\textwidth}{
    >{\raggedright\arraybackslash}p{0.26\textwidth}
    >{\raggedright\arraybackslash}X
}
\toprule
\textbf{Preguntas / palabras clave} &
\textbf{Notas / respuestas} \\
\midrule

\textbf{¿Qué problema práctico muestra el caso del demandado mayor?} &
Una cuota fijada a quien no entiende qué son los «alimentos» ni maneja bancos. Hay que notificarlo bien, explicarle y definir una forma de pago acreditable. \\

\textbf{¿Qué valor tiene una denuncia de violencia sin impulso?} &
Casi ninguno para la estrategia: la causa queda paralizada. Importa el estado real de cada causa, no sólo que exista la denuncia. \\

\textbf{¿Qué se hace si la mediación no puede notificar al requerido?} &
Si sólo hay domicilio laboral y el correo dice «domicilio desconocido», se cierra la mediación y se reabre si luego aparece un domicilio real. \\

\textbf{¿Cómo se actúa ante el incumplimiento de provisorios?} &
Se inicia o continúa la demanda, se intenta ejecutar los montos adeudados y, con tres cuotas impagas, puede pedirse la inscripción en el registro de deudores alimentarios. \\

\textbf{¿Qué estrategia puede usar el obligado que no paga?} &
Alegar que «no tiene» para presionar a la otra parte a aceptar menos de lo que se debe. El equipo decidió litigar hasta el final. \\

\midrule

\multicolumn{2}{p{\dimexpr\textwidth-4\tabcolsep\relax}}{
\textbf{Resumen:}
Los casos presentaron problemas prácticos: alimentos fijados a quien no los comprende, denuncias paralizadas por falta de impulso, mediación cerrada por domicilio desconocido y incumplimiento de provisorios, frente al cual se decidió ejecutar y litigar hasta el final.
} \\

\bottomrule
\end{tabularx}
\end{table}

% ============================================================
\chapter[Organización del equipo]{Bloque III · Organización del equipo y preparación de la reunión}
% ============================================================

\section{Revisión de escritos y fojas}

La revisión de los escritos antes de la reunión incluyó observaciones puntuales:

\begin{itemize}
    \item ubicar el dato «expediente» en el escrito;
    \item una de las causas va por fuera de la de violencia;
    \item un informe corresponde a otra hija; el informe del centro puede consultarse en el PJN;
    \item sobre el petitorio, se propuso hacerlo firmar, aunque la consultante querrá leerlo completo;
    \item respecto de las fojas, se mencionó la 122 y la 125; falta la que está en el medio.
\end{itemize}

% TODO: contenido incompleto o ambiguo en las notas originales (esta parte de la grabación es fragmentaria y las referencias a quién habla no son claras).

\section{Calendario de reuniones del equipo}

El equipo acordó reunirse el miércoles a las 18:30, porque la consultante se conecta a las 19 y no puede quedarse a las 20. Según el docente, para ese miércoles todo debe estar resuelto: el tiempo depende de los estudiantes, que deben tener el drive actualizado y las cosas listas. A modo de referencia, el docente indicó que la reunión anterior con ellos duró aproximadamente una hora y cuarto.

\section{Qué hay que decirle a la consultante antes de la reunión}

Antes de la reunión corresponde informar a la consultante sobre:

\begin{itemize}
    \item que se le fijaron \textbf{alimentos provisorios} que debe pagar, y explicarle la resolución;
    \item el tema del \textbf{gato}, que se le va a pedir;
    \item la contestación de la abogada contraria, que es extensa; se propuso que la consultante firme la última versión y que se lea con calma, con todos los datos ya chequeados.
\end{itemize}

\begin{important}
\textbf{Checklist previo a la reunión con la consultante:} (1) explicar la \textbf{resolución de alimentos provisorios} y su plazo de pago; (2) abordar el tema de la \textbf{mascota}; (3) recordar que \textbf{la comunicación con la abogada de la otra parte} corresponde sólo al patrocinio; (4) tener \textbf{firmados} la demanda y el petitorio.
\end{important}

\subsection*{Cuadro Cornell del Bloque III}

\begin{table}[H]
\centering
\begin{tabularx}{\textwidth}{
    >{\raggedright\arraybackslash}p{0.26\textwidth}
    >{\raggedright\arraybackslash}X
}
\toprule
\textbf{Preguntas / palabras clave} &
\textbf{Notas / respuestas} \\
\midrule

\textbf{¿Cómo se controla un escrito antes de firmarlo?} &
Se revisan los datos del expediente, las fojas citadas y las causas vinculadas, y se prevé que la consultante lea el petitorio completo. \\

\textbf{¿Cuándo y cómo se coordina el trabajo del equipo?} &
Mediante una reunión el miércoles a las 18:30, con el drive actualizado y todo resuelto de antemano; el tiempo de la reunión depende de la preparación. \\

\textbf{¿Qué hay que decirle a la consultante antes de la reunión?} &
Explicar la resolución de alimentos provisorios y su plazo, abordar el tema de la mascota, recordar que la comunicación con la abogada contraria corresponde sólo al patrocinio y tener firmados la demanda y el petitorio. \\

\midrule

\multicolumn{2}{p{\dimexpr\textwidth-4\tabcolsep\relax}}{
\textbf{Resumen:}
La organización del equipo exige revisar los escritos y las fojas, coordinar una reunión con el drive actualizado y preparar a la consultante sobre los alimentos provisorios, la mascota y los límites de comunicación.
} \\

\bottomrule
\end{tabularx}
\end{table}

% ============================================================
\chapter[Reunión con la consultante]{Bloque IV · Reunión con la consultante}
% ============================================================

\section{Documentación que lleva el docente}

El docente pidió imprimir el escrito, la demanda y la \textbf{resolución del provisorio}. La demanda ya fue contestada y el plazo vencía el jueves, por lo que las firmas debían realizarse ese día. También mencionó:

\begin{itemize}
    \item los \textbf{certificados de las atenciones}: hay tres turnos programados, dos más programados para mañana (martes) y el 13, y los próximos tres martes ya tienen turnos. Resulta conveniente declararlo, para dejar respaldado que \textbf{la atención continúa}; ya fue enviado por mail;
    \item los \textbf{informes de la psicóloga} de la niña: están en el expediente, porque fueron subidos y pedidos por el juzgado, por lo que se consideran útiles;
    \item los \textbf{estudios médicos}: no son necesarios.
\end{itemize}

\section{Los alimentos provisorios: obligación y estrategia}

\subsection{Contenido de la resolución}

Se fijaron alimentos provisorios para que la consultante pague al padre de la niña \textbf{mientras la niña viva con él}; por eso son provisorios. El monto es de \textbf{150.000 pesos}. La consultante fue notificada el 30 de septiembre y debe depositar \textbf{del 1 al 10 de octubre}. El docente explicó que no conoce sus posibilidades económicas y que no sabe si puede pagar.

La consultante manifestó su indignación: tardó siete años en lograr todo y el padre, en tres meses, obtuvo que le sacaran a la niña y los alimentos. El docente respondió que entiende su enojo, más allá de lo injusto que pueda resultar, pero que \textbf{está obligada al pago}: el equipo hará el escrito, pero la obligación rige.

\begin{important}
Ante una cuota provisoria que no se puede pagar completa: \textbf{depositar o transferir lo que se pueda} para \textbf{acreditar voluntad de pago}. Mientras se tramita el escrito, la obligación igualmente rige.
\end{important}

\subsection{Fundamento de la decisión judicial}

El padre lo solicitó a la jueza, y ésta lo concedió porque entiende que su situación económica empeoró: perdió el trabajo (consta el telegrama de despido). La consultante objetó que su jefe es su mejor amigo y que trabaja en una sociedad de dos amigos, de los cuales él es empleado. El docente indicó que eso es lo que entiende la jueza y que el niño necesita comer; la medida es provisoria y lo que genera enojo es la manera.

La consultante expresó una reacción violenta de enojo. El docente le advirtió que no lo diga, porque \textbf{«te están examinando absolutamente todo: lo que hacés, lo que decís»}.

\section{Límites éticos: no contactar a la abogada ni a la contraparte}

El docente abordó un tema incómodo, ya hablado anteriormente: la consultante llamó a la abogada de la otra parte. Los puntos expuestos fueron:

\begin{itemize}
    \item El juzgado marca los tiempos, ni el docente ni la consultante.
    \item Si la consultante no acuerda con las pautas del equipo, puede buscar otro patrocinio gratuito u otro abogado de la matrícula.
    \item El docente no puede actuar como si fuera la madre de la niña; debe ver el caso desde afuera y no desde adentro.
    \item La llamada \textbf{no puede volver a ocurrir}: coloca al docente en una falta de ética y, si se repite, el docente renunciará al expediente. Se trata de la primera y única advertencia.
    \item La consultante no puede hablar con la abogada, con el padre ni con su hija sobre el caso. Cuenta con un canal: el equipo atiende lunes y jueves, y el mail es su forma de comunicación, que se tolera.
    \item La llamada se hizo para entregar el gato; la abogada contraria podría denunciarla. La consultante está en un proceso judicial en el que \textbf{la examinada es ella}.
\end{itemize}

\begin{important}
\textbf{Pauta ética de la clínica.} El abogado no puede actuar como si fuera el familiar de su asistido: debe mirar el caso «desde afuera» para preservar la profesionalidad. La comunicación con la otra parte se canaliza \textbf{entre abogados}. Un contacto directo del asistido con la abogada contraria puede poner al patrocinante en una falta de ética y, ante una reiteración, el docente renunciaría al patrocinio.
\end{important}

\subsection{La entrega de la mascota}

El gato debe estar con la niña y la consultante debe entregarlo. El equipo propuso entregarlo después de las 19, en el hospital o en el estudio jurídico de la abogada contraria; la otra parte aún no respondió y tiene \textbf{tres días} para indicar qué y cuándo. Mientras tanto, el equipo está a derecho y no debe nada en el juzgado.

Si el juzgado dispone que se entregue en la casa del amigo, deberán ir: lo que menos se debe hacer es enojar al juzgado. El objetivo es que la niña y el juzgado vean que la niña tiene un lugar seguro y que se cumple lo que se dice.

Sobre la comunicación, la consultante puede estar en copia de los mails o manejarse con su madre; el equipo trata la mayoría de los asuntos presencialmente. La demanda debe ser firmada porque debe presentarse; los alimentos provisorios, la consultante debe pensarlos y avisar qué decidió.

\section{El CIF: cómo prepararse para la entrevista pericial}

\begin{definition}{CIF}
\textbf{Cuerpo Interdisciplinario Forense}: equipo de peritos que evalúa a las partes y a la niña por separado y produce un \textbf{informe} para el juzgado en cuestiones de familia, cuando el juzgado necesita conocer ciertas circunstancias. La entrevista es con un psicólogo, un psiquiatra y quizás un trabajador social (uno de los estudiantes no estaba seguro de la participación del trabajador social). Según la clase, de ese informe depende en gran medida la recomendación sobre la revinculación.
\end{definition}

Habrá una entrevista con la consultante, otra con la niña y otra con el padre, \textbf{por separado}, y luego el CIF elaborará un informe para el juzgado. Se señaló que ese informe tiene que salir bien para que la niña vuelva con la madre.

\subsection{Recomendaciones del docente para la entrevista}

\begin{enumerate}
    \item \textbf{Reconocer} lo que uno hizo mal y lo que podría haber hecho mejor.
    \item \textbf{Poner el foco en uno mismo y no en el otro progenitor}. Hablar de todo lo mal que está el padre «como padre, como persona y como individuo en el mundo» no sirve.
    \item \textbf{Hacer hincapié en la terapia}, mencionándola de buena manera. El docente aconsejó acostumbrarse a no decirlo mal, porque «se te va a escapar».
    \item \textbf{No atribuir soborno, complot o mala fe} a la Defensoría, a la psicóloga, al padre u otras instituciones: si la consultante habla mal de todos, el informe del CIF «va a dar mal».
    \item \textbf{Esperar preguntas incómodas e hirientes}: no es una entrevista con una amiga.
    \item No creer que todo el mundo está en contra.
\end{enumerate}

Sobre la consultante, que manifestó que en la Defensoría de Menores no la trataron bien, el docente le dijo que tiene la sensación de que ella cree que todo el mundo está en contra y que eso no la lleva a ningún lado; es una opinión personal y no técnica. La consultante respondió que no lo va a decir en el CIF pero que lo sigue pensando.

\subsection{El motivo del impedimento de contacto}

Un estudiante preguntó bajo qué motivo se impidió el contacto y se retiró a la niña (en el drive hay un archivo de 700 hojas) y si hay diagnóstico psiquiátrico. El docente respondió que existe un \textbf{informe reservado} que el equipo no vio y que ya se pidió ver. En principio tiene que ver con el comportamiento de la consultante en el hospital cuando fue a revisar a la niña, o con algo que dijo la niña; no se sabe qué ocurrió con exactitud, pero hay una pista sobre ese comportamiento. Ese hecho dio origen a otro informe que tampoco se vio.

La consultante negó la veracidad de lo relatado y sostuvo que una profesional nunca los atendió. El docente indicó que la mejor forma de que se aclare es que esas personas se presenten y declaren. La consultante mencionó firmas del equipo de orientación escolar, que no son del colegio sino del Ministerio de Educación, y que una señora la llamó a las 19. El docente explicó que es normal que se manejen en esos horarios, porque están sobrepasados; si no la llaman, hacen el informe indicando «no la encontré, no me atendió». Las entrevistas a la niña pueden hacerse sin consentimiento de la madre porque deben averiguar cómo es ella con la niña. En el caso de la hija de 31 años no hubo duda respecto de la consultante; en este caso las circunstancias dieron una duda que activó a las instituciones.

\begin{important}
\textbf{Recomendaciones para la entrevista del CIF:} (1) reconocer lo propio que pudo hacerse mejor; (2) poner el foco en uno mismo y no en el otro progenitor; (3) mencionar de forma positiva que se hace terapia; (4) no atribuir soborno, complot o mala fe a instituciones o peritos; (5) esperar preguntas incómodas; (6) no tomar al sistema como enemigo.
\end{important}

\section{Canales de comunicación y cierre de la reunión}

El docente reiteró que el equipo está presente los lunes y jueves, y que cualquier consulta se hace por mail a uno de los estudiantes. Insistió en que no todos están en contra de la consultante y que mucha gente trabaja a favor de la niña, \textbf{«que es la única que nos importa»}.

Sobre el gato: se espera que la otra parte responda y la consultante no habla con nadie, salvo con el docente. La consultante reconoció que quizá actuaron de manera imprudente por la desesperación. El docente recordó que \textbf{los abogados tienen que hablar entre ellos}. En caso de emergencia médica, se va al médico; si es una emergencia, por mail, porque no hay otro canal: la facultad no permite enviar WhatsApp y no se dispone del WhatsApp de nadie.

\subsection{Documentar la entrega de la mascota}

La consultante señaló que el escrito sobre la mascota no indica en qué domicilio estará ni a quién se entrega, y preguntó si podía agregarlo ella y hacerlo firmar a quien se lleva el gato. El docente respondió que no: el equipo ofrece \textbf{opciones} para que la otra parte conteste cuál elige y cuándo. Aun así, es conveniente que quede documentado, porque si el gato vuelve o no, la otra parte tendrá que responder.

\subsection*{Cuadro Cornell del Bloque IV}

\begin{longtable}{
    >{\raggedright\arraybackslash}p{0.26\textwidth}
    >{\raggedright\arraybackslash}p{0.68\textwidth}
}
\toprule
\textbf{Preguntas / palabras clave} &
\textbf{Notas / respuestas} \\
\midrule
\endhead

\textbf{¿Qué documentación se entrega y se firma?} &
La demanda, el petitorio y la resolución del provisorio; los certificados de atención, que dejan respaldado que la atención continúa; y los informes de la psicóloga de la niña, que ya están en el expediente. Los estudios médicos no son necesarios. \\

\textbf{¿Qué hacer si no se puede pagar la cuota provisoria completa?} &
Depositar o transferir lo que se pueda para acreditar voluntad de pago; la obligación rige igualmente mientras se tramita el escrito. \\

\textbf{¿Por qué no se puede llamar a la abogada de la otra parte?} &
Compromete la ética del patrocinante, puede dañar la estrategia y, si se repite, el docente renunciaría al patrocinio. La comunicación es entre abogados. \\

\textbf{¿Qué es el CIF?} &
Cuerpo Interdisciplinario Forense: peritos (psicólogo, psiquiatra y quizás trabajador social) que entrevistan por separado a las partes y a la niña y elaboran un informe para el juzgado. \\

\textbf{¿Cómo prepararse para la entrevista del CIF?} &
Reconocer lo que pudo hacerse mejor, enfocarse en uno mismo, valorar la terapia, no hablar mal del otro progenitor ni de instituciones y esperar preguntas incómodas. \\

\textbf{¿Por qué documentar la entrega de la mascota?} &
Porque debe constar a quién y dónde se entrega; así queda constancia si la otra parte responde o no. El equipo ofrece opciones y la otra parte debe indicar cuál elige y cuándo. \\

\textbf{¿Por qué canal se comunica el asistido?} &
Con el equipo, presencialmente (lunes y jueves) o por mail; en emergencia, por mail, porque no hay otro canal. La facultad no permite WhatsApp. \\

\midrule

\multicolumn{2}{p{\dimexpr\textwidth-4\tabcolsep\relax}}{
\textbf{Resumen:}
En la reunión el docente combinó claridad técnica (la cuota provisoria rige y debe acreditarse voluntad de pago) con límites éticos (no contactar a la abogada ni a la contraparte, bajo riesgo de renuncia al patrocinio) y una preparación estratégica para el CIF: reconocer lo propio, enfocarse en uno mismo y no atribuir mala fe a las instituciones.
} \\

\bottomrule
\end{longtable}


% ============================================================
\chapter[Familia ampliada y cierre]{Bloque V · Intervención de la hermana y cierre}
% ============================================================

\section{Revisión del escrito por la hermana}

La hermana de la consultante, que estudió un año de Derecho antes de dejar, revisó el escrito y realizó correcciones respecto de su nombre y del vínculo con su hija: el escrito decía «prima», pero su hija es sobrina de la niña, porque la niña es su hermana. Además, en una parte el escrito indica que la niña tiene ocho años y en otra siete. Leyó rápido, fijándose en nombres y vínculos.

\begin{important}
\textbf{Lección de redacción.} Los datos de vínculo y edad deben ser coherentes en todo el escrito (aparecieron inconsistencias: «prima» por «sobrina», ocho años y siete años para la misma niña). Un escrito se revisa \textbf{con el cliente} antes de firmar.
\end{important}

\section{Régimen de contacto y mensajes}

La hermana explicó que el cuidado era compartido porque no había motivo para que el padre no pudiera ver a la niña: de hecho, la veía todos los días en el colegio.

El docente se quedó con el documento para leerlo y anunció que habrá un nuevo informe, elaborado por el equipo. Si ese informe sale bien, hay muchas probabilidades de que se recomiende la \textbf{revinculación}.

\subsection{Representación de la hermana}

La hermana preguntó si necesitaba representación aparte o si podía actuar a través del equipo, para tener días con la niña, ya que no obtuvo respuesta del padre. El docente indicó que ya está incorporada en el expediente y que la representación puede darse por este equipo; la etapa es la preliminar.

\subsection{Mensajes enviados al padre}

La hermana relató que, cuando el padre solicitó al juzgado que ella dejara de enviar mensajes, las capturas que adjuntó mostraban sólo los mensajes en que él no contestaba, mientras que las primeras dos veces sí le contestó, incluso de forma amenazante. Lo que le llegó no es una restricción: se pide que cese el envío de mensajes hasta que esté la evaluación. Esa evaluación ya la tuvo y ahora se aceptó hacer la evaluación de todos de nuevo.

% TODO: contenido incompleto o ambiguo en las notas originales (la grabación termina mientras la hermana relataba lo que el padre habría dicho a otra hermana, la madrina; y la última frase del docente es fragmentaria).

\begin{note}
Según el docente, cuando los estudiantes ingresaron al caso ya era tarde para hacer otra cosa respecto de lo ya ocurrido.
\end{note}

\begin{important}
\textbf{Conclusión general de la clase.} La práctica del patrocinio exige \textbf{informes completos}, \textbf{comunicación disciplinada} (solo entre abogados), \textbf{honestidad con el asistido} sobre lo provisorio y lo obligatorio, y una \textbf{preparación estratégica} de las entrevistas periciales. El equipo funciona si cada integrante entrega la información que sólo él conoce.
\end{important}

\subsection*{Cuadro Cornell del Bloque V}

\begin{table}[H]
\centering
\begin{tabularx}{\textwidth}{
    >{\raggedright\arraybackslash}p{0.26\textwidth}
    >{\raggedright\arraybackslash}X
}
\toprule
\textbf{Preguntas / palabras clave} &
\textbf{Notas / respuestas} \\
\midrule

\textbf{¿Por qué se revisa el escrito con la familia?} &
Para detectar errores de vínculo y edad (por ejemplo, «prima» por «sobrina», ocho y siete años) antes de presentarlo. \\

\textbf{¿Puede la hermana participar y cómo?} &
Ya está incorporada en el expediente y su representación puede darse por el mismo equipo. \\

\textbf{¿Qué ocurre si el nuevo informe sale bien?} &
Hay muchas probabilidades de que se recomiende la revinculación. \\

\midrule

\multicolumn{2}{p{\dimexpr\textwidth-4\tabcolsep\relax}}{
\textbf{Resumen:}
La revisión del escrito con la familia permite detectar inconsistencias de vínculo y edad. La hermana está incorporada al expediente y puede ser representada por el equipo; un nuevo informe favorable aumentaría las probabilidades de que se recomiende la revinculación.
} \\

\bottomrule
\end{tabularx}
\end{table}

% ============================================================
\chapter[Síntesis final]{Síntesis final de la clase}
% ============================================================

\begin{table}[H]
\centering
\begin{tabularx}{\textwidth}{
    >{\raggedright\arraybackslash}p{0.26\textwidth}
    >{\raggedright\arraybackslash}X
}
\toprule
\textbf{Preguntas / palabras clave} &
\textbf{Notas / respuestas} \\
\midrule

\textbf{¿Qué datos mínimos debe tener todo informe de caso?} &
Juzgado completo, medidas provisorias con su vencimiento, deuda y pagos, grupo familiar, situación laboral de las partes y respuestas de los oficios. \\

\textbf{¿Cuál es la prueba de que un informe está mal hecho?} &
Que haya que ir al expediente para entender el caso. \\

\textbf{¿Cómo se actúa ante el incumplimiento de provisorios?} &
Se ejecutan los montos adeudados y, con tres cuotas impagas, puede pedirse la inscripción en el registro de deudores alimentarios; el equipo decidió litigar hasta el final. \\

\textbf{¿Qué debe hacer la consultante si no puede pagar la cuota provisoria completa?} &
Depositar o transferir lo que pueda para acreditar voluntad de pago. \\

\textbf{¿Por qué no se puede llamar a la abogada de la otra parte?} &
Compromete la ética del patrocinante y, si se repite, el docente renunciaría al patrocinio. \\

\textbf{¿Cómo prepararse para el CIF?} &
Reconocer lo propio, enfocarse en uno mismo, valorar la terapia, no atribuir mala fe a las instituciones y esperar preguntas incómodas. \\

\midrule

\multicolumn{2}{p{\dimexpr\textwidth-4\tabcolsep\relax}}{
\textbf{Resumen:}
La clase enseñó que el patrocinio jurídico en familia se apoya en dos pilares: un \textbf{informe de caso completo}, que cualquier colega pueda leer sin ir al expediente, y una \textbf{relación disciplinada con el asistido}. En la práctica aparecieron problemas reales (alimentos fijados a quien no entiende la obligación, denuncias paralizadas, mediación cerrada por domicilio desconocido e incumplimiento de provisorios). En la reunión con la consultante, el docente combinó claridad técnica con límites éticos y una preparación estratégica para el CIF. El hilo conductor es la \textbf{profesionalidad}: mirar el caso desde afuera, trabajar en equipo y documentar todo.
} \\

\bottomrule
\end{tabularx}
\end{table}

% Nota editorial: el apunte original incluye un recuadro de «Marco normativo de referencia» que
% se declara complementario y no citado en la clase; se omitió por no formar parte del contenido
% expuesto por el docente.

\end{document}
